% LSPV-3 REPORT::ORBITBATCH1::2026-10-05T02:35:00-07:00
\documentclass[11pt]{article}
\usepackage[margin=.86in]{geometry}
\usepackage{amsmath,amssymb}
\begin{document}
\section*{LSPV-3 ORBITBATCH1: orbit-query deduplication and normalization audit}
\textbf{GLOBAL STATUS: OPEN.} P versus NP remains open. This report proves source-side query-set identities and conservative accounting for the declared SU05/Hadamard realization. It proves no SAT shrinker, no arbitrary-algorithm acquisition theorem, and no forbidden-region contradiction.

\section*{CORPUS REVIEWED}
Repository governance, CODEX, the pnp-research skill and reconstruction mode, root main.tex, TASK-001, newest verified MAIN surfaced by repository search, LSPV-2 SHORTKERN1, LSPV-3 STREAMCOMP1/SOURCEEDGE1/SOURCETRACE1/LOCALBIT1/SEEDBRANCH1/EXTEVAL1, LSPV-4 BATCHSETUP1/CLASSCHARGE1, and the HEC/CRS-PPC/CCSPC acquisition lineage were reviewed. Repository-wide search was used for prior LDE acquisition, setup sharing, classification, and multipoint-evaluation claims. The immediate LSPV-2 predecessor remains SHORTKERN1, which is substantive and has its verification sentinel.

\section*{Target and exact quantifiers}
Fix one declared source-sanctioned SU05/Hadamard realization, one legal RMV parent polynomial, and one legal MV line restriction
\[
f:B=\mathbb F_p^d\to\mathbb F_p,\qquad d=r/2.
\]
For every orbit index \(i\) used by the extractor and every seed \(v\in B\), the declared q-ary query is \(f(A^iv)\), where \(A\in GL_d(\mathbb F_p)\). The target is to charge the union of these query addresses over all seeds without confusing semantic seed multiplicity with distinct oracle evaluations.

\section*{PROVED/EXACT: orbit-query deduplication}
For every fixed \(i\), invertibility of \(A^i\) implies
\[
\{A^iv:v\in B\}=B.
\]
Hence for any nonempty set \(I\) of orbit indices,
\[
\boxed{\bigcup_{i\in I}\{A^iv:v\in B\}=B.}
\]
Therefore all q-ary field evaluations needed over every vector seed \(v\) and every declared orbit position can be serviced from the complete evaluation table of \(f\) on \(B\), containing exactly
\[
|B|=p^d=p^{r/2}
\]
field symbols. Multiplying the vector-seed count by the number of orbit positions is an overcount of distinct oracle addresses.

This is an address-set identity only. It does not collapse the number of emitted extractor strings, Hadamard coordinates, or consumers.

\section*{Direct-output volume survives}
The declared Hadamard seed contains the vector seed plus a binary-code coordinate of \(\log_2p+O(1)\) bits. Thus the number of seed presentations per line remains
\[
p^d\cdot p^{1+o(1)}=p^{r/2+1+o(1)}.
\]
There are \(2p^{r/2}\) distinct MV lines. Consequently, if RMV explicitly emits one direct extractor output for every line and every full seed, the direct output count alone is
\[
2p^{r+1+o(1)}.
\]
With the source-side normalization \(N=p^r\), this is
\[
\boxed{N^{1+1/r+o(1)}}.
\]
Thus orbit-query deduplication can reduce field-evaluation work, but cannot reduce explicit output-writing below the number of emitted direct outputs.

\section*{Shared parent and line setup}
For a parent oracle \(F:B^2\to\mathbb F_p\), the two MV line families are
\[
f_{1,b}(z)=F(z,b),\qquad f_{2,b}(z)=F(b,z).
\]
All children are restrictions of one parent realization. No source theorem requires independently reacquiring a fresh parent LDE for each child. In BATCHSETUP1 terminology, one already-acquired parent evaluator can be a single reusable setup state, so the source-side candidate has \(\delta_{\rm parent}=0\). This is conditional on retaining or recomputing a SAME\_REALIZATION parent evaluator with its time and space fully charged. It is not an arbitrary-SAT acquisition theorem.

CLASSCHARGE1 corrects the generic ledger: total classification exponent \(\chi\) is a separate max term. The conservative edge accounting is
\[
C_i\ge\max\{r_i,\chi,\delta+p,\eta+g+\theta C_j\},
\]
not a generic sum \(\delta+p+\kappa\).

\section*{Primary-literature check: batched polynomial evaluation}
Bhargava, Ghosh, Guo, Kumar, and Umans give deterministic multivariate multipoint evaluation over every finite field in
\[
(d_0^m+Q)^{1+o(1)}\operatorname{poly}(m,d_0,\log|\mathbb F|)
\]
time for an \(m\)-variate polynomial of individual degree below \(d_0\), given by coefficients, at \(Q\) points. This is a directly relevant algebraic upper-bound tool after coefficients are already available. It reinforces that repeated pointwise Horner evaluation is not a mandatory charge. It does not provide coefficient acquisition from an arbitrary SAT computation.

For the present orbit query set, an even simpler source-specific fact dominates: the distinct requested points are exactly \(B\). If the complete child table is materialized, its output volume is already \(p^d\), so no asymptotic field-evaluation lower bound larger than that may be inferred merely from seed-by-orbit multiplicity.

\section*{Normalization audit: decisive caveat}
The normalization
\[
N=p^r
\]
is the cardinality of the full field evaluation domain. It is not the original SAT input length and not the vMS PCPP proof length. In the checked vMS instantiation,
\[
h=m^{100},\qquad r=\min\{2^j:h^r\ge |y|\},\qquad
\Delta=hr,\qquad p\in[\Delta^{100},2\Delta^{100}].
\]
Therefore
\[
p^r\ge (hr)^{100r},
\]
whereas the canonical proof table occupies only \(h^r\) grid positions before extension. The ratio is enormous:
\[
\frac{p^r}{h^r}\ge r^{100r}h^{99r}.
\]
Accordingly, a source-side exponent such as \(N^{1.25}\) or \(N^{1.5}\) has no direct meaning as a SAT exponent. Any forbidden-region comparison must first derive an explicit map from original SAT length \(n\) to every source length, proof length, field-domain size, and returned SAT-instance length. The earlier source-side fixed points are diagnostic only.

\section*{Falsification}
\textbf{FALSIFIED/RETIRED:} charging \(m\) independent field evaluations for every vector seed when counting distinct oracle addresses. Orbit invertibility makes the union of addresses exactly the child domain.

\textbf{FALSIFIED/RETIRED:} interpreting contraction relative to \(N=p^r\) as a SAT-input contraction without a source-to-SAT length map.

\textbf{NOT FALSIFIED:} explicit output volume \(N^{1+1/r+o(1)}\) for the declared all-line/all-seed direct branch; parent-evaluator acquisition; retained-space cost; SAME\_REALIZATION return; arbitrary-algorithm acquisition.

\section*{Hostile controls and barriers}
RECOMP may discard the shared parent and pay repeated acquisition; DELAY may materialize only queried lines; INTERP may use a compact coefficient or tensor-grid evaluator; COPY changes presentations but not distinct addresses; semantic and algebraic oracles can make parent evaluation free. Horn, 2-SAT, affine/Tseitin, bounded-treewidth, and independent-product controls remain relevant to any claimed hardness transfer. The address identity and multipoint-evaluation upper bound are algebraic and relativizing. No nonrelativizing or nonalgebrizing hinge is supplied.

\section*{DUPLICATION / CONTRADICTION AUDIT}
EXTEVAL1 correctly separated representation acquisition from local evaluation but used per-query evaluation accounting. ORBITBATCH1 sharpens it by deduplicating the complete declared orbit query set. BATCHSETUP1 already allowed globally shared setup; ORBITBATCH1 identifies the common parent restriction structure that can realize such sharing source-side. CLASSCHARGE1 supersedes the root's generic \(\delta+p+\kappa\) shorthand with separate total classification exponent \(\chi\). No contradiction with STREAMCOMP1 occurs because STREAMCOMP1 is a conservative simulation ledger, not a lower bound requiring independent generation.

\section*{WHAT CHANGED RELATIVE TO ROOT MAIN}
The root LSPV-3 region is stale in two respects: it does not yet include SEEDBRANCH1/EXTEVAL1, and it retains \(\delta+p+\kappa\) as a conservative candidate despite verified CLASSCHARGE1 showing that generic sequential classification must be represented separately. The root also needs the new normalization guard: \(p^r\) is a field-domain size, not SAT input length.

\section*{ADOPTED / REJECTED / RETIRED PRIOR COMPONENTS}
Adopt STREAMCOMP1, BATCHSETUP1, CLASSCHARGE1, SOURCETRACE1, LOCALBIT1, SEEDBRANCH1, EXTEVAL1 at their declared scopes. Reject independent reacquisition of every line as source-forced. Retire seed-by-orbit query multiplicity as a distinct-address count. Retire any direct use of \(N=p^r\) fixed points in SAT forbidden-region comparisons.

\section*{STALE OR OVERSTATED CLAIMS FOUND}
The root formula \(\delta+p+\kappa\) is no longer the preferred implementation-independent classification ledger after CLASSCHARGE1. Also, EXTEVAL1's favorable source-side exponents remain valid only in the field-domain normalization and must not be described as SAT exponents.

\section*{Single next target}
Derive the complete length map
\[
n_{\rm SAT}\longmapsto |y|\longmapsto (h,r,\Delta,p)\longmapsto p^r
\]
for the exact vMS PCPP branch, including the polynomial exponent hidden in proof length as a function of the original SAT instance. Then compose ORBITBATCH1, output-writing volume, parent acquisition, retained space, and any SAT return length in the original \(n_{\rm SAT}\) normalization. This is required before any Fortnow--Lipton--van Melkebeek--Viglas or Buss--Williams forbidden-region comparison.

\section*{References}
V. Bhargava, S. Ghosh, Z. Guo, M. Kumar, and C. Umans, \emph{Fast Multivariate Multipoint Evaluation Over All Finite Fields}, FOCS 2022; ECCC TR22-063.

R. Shaltiel and C. Umans, \emph{Simple Extractors for All Min-Entropies and a New Pseudorandom Generator}, J. ACM 52(2):172--216, 2005.

D. van Melkebeek and N. Mocelin Sdroievski, \emph{Instance-Wise Hardness and Refutation versus Derandomization for Arthur--Merlin Protocols}, Computational Complexity 34:15, 2025.

\bigskip\noindent\textbf{VERIFICATION SENTINEL: ORBITBATCH1 complete; P versus NP OPEN.}
% END_OF_VERIFIED_REPORT::LSPV-3::ORBITBATCH1
\end{document}
